
\documentclass{article}
\usepackage{colortbl}
\usepackage{makecell}
\usepackage{multirow}
\usepackage{supertabular}

\begin{document}

\newcounter{utterance}

\centering \large Interaction Transcript for game `function\_detective', experiment `two\_numbers\_passive\_labeled\_pairs\_oneshot', episode 299 with Qwen3.6{-}35B{-}A3B{-}FP8{-}without{-}reasoning.
\vspace{24pt}

{ \footnotesize  \setcounter{utterance}{1}
\setlength{\tabcolsep}{0pt}
\begin{supertabular}{c@{$\;$}|p{.15\linewidth}@{}p{.15\linewidth}p{.15\linewidth}p{.15\linewidth}p{.15\linewidth}p{.15\linewidth}}
    \# & \multicolumn{2}{c}{Player} && \multicolumn{2}{c}{Game Master} \\
    \hline

    \theutterance \stepcounter{utterance}  
    & & & \multicolumn{4}{p{0.6\linewidth}}{
        \cellcolor[rgb]{0.9,0.9,0.9}{
            \makecell[{{p{\linewidth}}}]{
                \texttt{\tiny{[P1$\langle$GM]}}
                \texttt{I have a hidden Python function for you to figure out.} \\
\\ 
\texttt{Your job is to infer what it does and then submit a correct Python implementation.} \\
\\ 
\texttt{Function signature:} \\
\texttt{(a: int, b: int) {-}> int} \\
\\ 
\texttt{You have 1 turns total.} \\
\\ 
\texttt{I have already given you a fixed list of candidate input{-}output pairs with true or false labels. Solve using those statements only.} \\
\\ 
\texttt{You must make your best SOLVE attempt no later than the second{-}to{-}last turn.} \\
\texttt{Do not wait until the very end.} \\
\\ 
\texttt{You can use these actions:} \\
\texttt{{-} JSON format described at the end of the prompt} \\
\\ 
\\ 
\texttt{Do not add import statements in your guess.} \\
\texttt{Anything you might need is already available, so just write the function body in hypothesis and return the answer.} \\
\\ 
\texttt{Preloaded examples:} \\
\texttt{1. (((1, 10), 3), True)} \\
\texttt{2. ((({-}10, 1), 14), False)} \\
\texttt{3. (((10, {-}99), 4), True)} \\
\texttt{4. ((({-}99, 0), 1), False)} \\
\texttt{5. ((({-}1000, 0), 0), True)} \\
\texttt{6. ((({-}2, {-}1), 1), False)} \\
\texttt{7. (((99, 0), 0), True)} \\
\texttt{8. ((({-}1000, {-}99), 7), False)} \\
\texttt{9. ((({-}99, 2), 5), True)} \\
\texttt{10. (((1000, {-}2), 1), False)} \\
\\ 
\texttt{Scoring:} \\
\texttt{{-} correctness is required for any non{-}zero score} \\
\texttt{{-} in interactive modes, solving with fewer information requests is better} \\
\\ 
\texttt{Input formatting reminders:} \\
\texttt{{-} strings must be quoted} \\
\texttt{{-} booleans must be True/False} \\
\texttt{{-} use commas between multiple values} \\
\texttt{{-} do not write any extra action tags in the same message} \\
\\ 
\texttt{Format example:} \\
\texttt{You write:} \\
\texttt{\{'current\_mode': <probing/confirming/don't know; choose whichever best describes your current state in the task>, 'input\_rationale': <brief description of why you chose this hypothesis>, 'hypothesis': <either the exact string "unknown" if you have no guess yet OR your current best guess as runnable Python matching the given signature, with a one{-}line docstring stating what it does. Never write placeholder code, if you cannot commit to a guess, use "unknown". Match the given signature.} \\
\texttt{e.g.} \\
\texttt{def f(x):} \\
\texttt{	"""Returns the absolute value of x."""} \\
\texttt{	return x if x >= 0 else {-}x} \\
\\ 
\texttt{Use "unknown" if you have no guess.>, 'confidence\_score': <score between 1{-}100. 100 means highest confidence in your hypothesis>\}} \\
\\ 
\texttt{Output only one valid json per turn.} \\
            }
        }
    }
     \\ \\

    \theutterance \stepcounter{utterance}  
    & \multicolumn{4}{p{0.6\linewidth}}{
        \cellcolor[rgb]{0.561,0.737,0.561}{
            \makecell[{{p{\linewidth}}}]{
                \texttt{\tiny{[P1$\rangle$GM]}}
                \texttt{Let's analyze the input{-}output pairs to infer the function's behavior.} \\
\\ 
\texttt{The function signature is `f(a: int, b: int) {-}> int`, but the examples show a tuple of three elements: `(a, b, c)` and a boolean result. Wait, looking closely at the examples:} \\
\texttt{`(((1, 10), 3), True)`} \\
\texttt{This looks like the input to the hidden function is actually a tuple `(a, b)` and the third element `c` might be part of the test case structure or the function takes 3 arguments?} \\
\texttt{No, the signature is `(a: int, b: int) {-}> int`.} \\
\texttt{However, the examples are formatted as `(((a, b), c), label)`.} \\
\texttt{This suggests the "input" to the test harness is `(a, b, c)` and the output is a boolean.} \\
\texttt{But the function signature says it returns an `int`.} \\
\texttt{Let's re{-}read carefully: "I have a hidden Python function... Function signature: (a: int, b: int) {-}> int".} \\
\texttt{"You have already given you a fixed list of candidate input{-}output pairs with true or false labels."} \\
\\ 
\texttt{This implies the "input{-}output pairs" are actually test cases for a predicate or a property check, OR the hidden function returns an integer, and the "True/False" label indicates whether the returned integer satisfies some condition, OR the hidden function is actually a predicate `(a,} \\
            }
        }
    }
    & & \\ \\

    \theutterance \stepcounter{utterance}  
    & & & \multicolumn{2}{p{0.3\linewidth}}{
        \cellcolor[rgb]{0.9,0.9,0.9}{
            \makecell[{{p{\linewidth}}}]{
                \texttt{\tiny{[GM$|$GM]}}
                \texttt{Game Verdict: LOSS (turn limit reached)} \\
            }
        }
    }
    & & \\ \\

\end{supertabular}
}

\end{document}
